\section{Conclusion and Future Work}
We studied selfless traffic routing in SUMO with a cooperative MAPPO policy: a
centralized critic, a shared per-candidate route actor, and a difference-reward team
objective in seconds of travel time, with a congestion-gated incentive that pays for a
detour only when it relieves a real jam. On a congested New York City corridor, deployed
greedily on 20 disjoint held-out scenarios, it completes every vehicle and cuts mean fleet
travel time by 10.4\% (15/20 wins, $p{=}0.02$), with a lower tail and fewer deadline
misses---executing genuine load-spreading detours nearly 40\% of the time, not by
replanning alone.

The more interesting finding is conditional. Ablating the guard so that even the pile-on
detours execute helps a further few percent, but almost all of it falls on the most
congested scenarios and vanishes in light traffic. Selfless routing pays in proportion to
the network's price of anarchy---so the gain is largest where congestion is worst (up to
$-27\%$ on the most congested block), and a penetration sweep shows it growing monotonically
with the controlled share, reliable once the policy routes at least a quarter of the fleet.

Several directions follow. Most directly motivated by our data is evaluation on
high-price-of-anarchy networks (Pigou- or Braess-style bottlenecks), where selfless routing
has a large, persistent margin rather than the intermittent one a dense grid offers. The
weak learning signal and the detour-rate drift invite stronger credit assignment (a
counterfactual critic~\citep{foerster2018coma} or marginal-cost pricing), and partial
adoption calls for training under mixed-autonomy traffic. Finally, a graph-neural-network
encoding would let a policy transfer across networks, toward a deployable controller.
